\section{Equality and the number of boundary components}
\label{n3:rigidity:setup}

We now classify equality on the original initial data. The numerical
inequality has already been proved for a larger boundary class; that
larger class is necessary for the variations below. We retain the
constraint, ADM, and spacetime conventions of
Equations~\eqref{n3:intro:constraints}--\eqref{n3:intro:k-sign}.
In particular $J$ is a covector, the dominant energy condition is
$\mu\geq|J|_g$, and $K=\tfrac12\mathcal L_n g$ for the future normal.
All data remain smooth; the asymptotic hypotheses require only the two
metric derivatives and one $K$ derivative in
Equation~\eqref{n3:intro:decay}, with $q>1/2$.

\paragraph{The numerical input and the cut convention.}
Theorem~\ref{n3:intro:exterior-inequality} applies to a connected orientable
smooth three-dimensional exterior, complete as a metric space with its
nonempty compact smooth boundary included, and with exactly one AF end.
For smooth data with $g-\delta=O_2(r^{-q})$,
$K=O_1(r^{-1-q})$, $q>1/2$, integrable $\mu$ and $|J|_g$, existing
finite ADM limits, $\mu\geq|J|_g$, $E>|P|_\delta$, and
$\theta_+(\partial\Omega)\leq0$ for the normal into the exterior, it gives
\[
\sqrt{E^2-|P|_\delta^2}\geq
\sqrt{\frac{a_g(\partial\Omega)}{16\pi}}.
\]
It requires no fill-in and no outermostness, outer area minimization,
or connectedness of the boundary. In the variational proof we apply
this established theorem to nearby data on the same one-sided exterior.

The infimum is over the enclosing cuts of
Definition~\ref{n3:intro:exterior-class}: each cut separates the entire
inner boundary from the distant part of the unique end, may be
disconnected, and may coincide with the original boundary. The normal
points toward that end, and coincident frontier area is counted.
Enclosing cuts and minimizing sets are taken with bounded complementary
pockets filled. During a perimeter argument we may enlarge the
competitor class to all relatively compact finite-perimeter sets
containing a fixed smooth inner obstacle. This class is closed under
unions and intersections; its minimizers fill bounded complementary
pockets, because filling such a pocket removes perimeter. Its infimum
agrees with the smooth-cut infimum by the outward approximation proved
in Section~\ref{n3:variation:area-subsection}.

\begin{definition}[Partial-coincidence outermostness]
\label{n3:rigidity:partial-outermostness}
Let the boundary of a one-ended exterior be
$S=\bigsqcup_{i=1}^{\ell}S_i$, with $1\leq\ell<\infty$ and each
$S_i$ connected. The boundary is outermost in the partial-coincidence
sense if no enclosing cut $\Gamma\ne S$ has all of the following
properties: every connected component of $\Gamma$ is either an entire
original component $S_i$ or is wholly contained in the open exterior;
at least one component lies in the open exterior; and
$\theta_+(\Gamma)\leq0$ everywhere for the normal toward the end.
\end{definition}

This condition includes comparisons that retain some original
components and replace others by interior components. It will be used
only after contact regularity has shown that the active minimizing
cuts have exactly this form. We make no claim here for the weaker
condition excluding only wholly interior enclosing cuts.

\begin{theorem}[Finite-component exterior equality rigidity]
\label{n3:intro:rigidity}
Let $(M,g,K)$ be a smooth connected orientable complete initial data
set without boundary, consisting of a compact core and finitely many
asymptotically flat ends. On each end assume
$g-\delta=O_2(r^{-q})$, $K=O_1(r^{-1-q})$, with $q>1/2$, and existing
finite ADM flux limits. Assume that $\mu$ and $|J|_g$ are integrable
and that $\mu\geq|J|_g$ everywhere.

Choose an end and a finite nonempty disjoint union
\[
S=\bigsqcup_{i=1}^{\ell}S_i,\qquad 1\leq\ell<\infty,
\]
of compact connected smooth embedded two-sided surfaces without boundary,
separating
that end from a complementary region. Let $\Mext$ be the closure of
the connected side toward the chosen end. Suppose that
$\partial\Mext=S$, that $\Mext$ has exactly this one end, and that,
for the normal pointing into $\Mext$:
\begin{enumerate}[label=\textup{(\roman*)}]
\item every $S_i$ is a future MOTS, so $H_{S_i}+\tr_{S_i}K=0$;
\item $S$ is outermost in the partial-coincidence sense of
Definition~\ref{n3:rigidity:partial-outermostness};
\item every enclosing cut in the closed exterior has area at least
$A=|S|_g=\sum_{i=1}^{\ell}|S_i|_g>0$.
\end{enumerate}
Let $(E,P)$ be the ADM vector of the chosen end. If
\begin{equation}
\label{n3:intro:equality}
E>|P|_\delta,\qquad
m=\sqrt{E^2-|P|_\delta^2}=\sqrt{\frac{A}{16\pi}},
\end{equation}
then $\ell=1$ and $S$ is diffeomorphic to a sphere. There is a proper
smooth spacelike embedding of the entire original $\Mext$ into the
maximally extended Schwarzschild spacetime of mass $m$, inducing $g$
and the given $K$ for a consistent future unit normal. The chosen end
approaches the corresponding spatial infinity. The image of $S$ is a
smooth cross-section of the corresponding future horizon or the
bifurcation sphere. The embedding and its future normal are smooth
up to $S$, and the image extends as a smooth spacelike hypersurface
through that boundary.
\end{theorem}

\begin{corollary}[Strictness for disconnected boundary]
\label{n3:rigidity:finite-strictness}
Under the hypotheses of Theorem~\ref{n3:intro:rigidity}, retain
$E>|P|_\delta$ but do not assume its equality in mass and area.
If $\ell\geq2$, then
\[
\sqrt{E^2-|P|_\delta^2}>\sqrt{\frac{A}{16\pi}}.
\]
\end{corollary}

\begin{proof}
The original exterior is complete as a metric space with boundary,
as follows. An intrinsic Cauchy sequence is Cauchy in the complete
ambient manifold, and its ambient limit remains in the closed exterior.
A smooth interior or boundary half-ball chart then gives convergence
in the intrinsic distance. Its compact
boundary and unique end place it in
Definition~\ref{n3:intro:exterior-class}. Since $S$ is an enclosing cut
and every enclosing cut has area at least $A$, its enclosing infimum
is $A$. Theorem~\ref{n3:intro:exterior-inequality} gives the non-strict
inequality. Equality would imply $\ell=1$ by
Theorem~\ref{n3:intro:rigidity}, contradicting $\ell\geq2$.
\end{proof}

The proof of Theorem~\ref{n3:intro:rigidity} occupies the next two
sections. Each $S_i$ is orientable because it is two-sided in the
orientable manifold $M$, but no genus is prescribed. The proof first
constructs a causal stationary field on the original exterior and
concentrates its area multiplier on the entire boundary. A single
complete conformal-double argument then gives positive intrinsic
curvature and the same area $A$ on every component, forcing
$\ell=1$. Schwarzschild reconstruction follows with connectedness
established by the argument. The conclusion concerns only the chosen
exterior; it allows a nonconstant time graph and nonzero momentum in
the original asymptotic frame.
